% paper96-copilot-cover-design.tex
% LLM-Guided Cover Design: Copilot Participation in Decomposition
% JuGeo Paper Series — Paper 96
% Compile: pdflatex paper96-copilot-cover-design

\documentclass[11pt]{article}
% jugeo-common.tex — shared preamble for all Judgment Geometry papers
% Include via: % jugeo-common.tex — shared preamble for all Judgment Geometry papers
% Include via: % jugeo-common.tex — shared preamble for all Judgment Geometry papers
% Include via: \input{jugeo-common}

% ── Packages ────────────────────────────────────────────────────────────
\usepackage{amsmath,amssymb,amsthm,mathtools}
\usepackage{tikz-cd}
\usepackage{listings}
\usepackage{xcolor}
\usepackage{xspace}
\usepackage{booktabs}
\usepackage{multirow}
\usepackage{enumitem}
\usepackage[expansion=false]{microtype}
\usepackage{hyperref}
\usepackage{cleveref}
% \usepackage{thmtools}  % disabled: not available in this TeX installation
\usepackage{float}
\usepackage{caption}
\usepackage{subcaption}
\usepackage{tabularx}
\usepackage{stmaryrd}       % \llbracket, \rrbracket
\usepackage{bbm}            % \mathbbm{1}

% ── Page geometry ───────────────────────────────────────────────────────
\usepackage[margin=1in]{geometry}

% ── Theorem environments ────────────────────────────────────────────────
\theoremstyle{plain}
\newtheorem{theorem}{Theorem}[section]
\newtheorem{lemma}[theorem]{Lemma}
\newtheorem{proposition}[theorem]{Proposition}
\newtheorem{corollary}[theorem]{Corollary}

\theoremstyle{definition}
\newtheorem{definition}[theorem]{Definition}
\newtheorem{example}[theorem]{Example}
\newtheorem{construction}[theorem]{Construction}

\theoremstyle{remark}
\newtheorem{remark}[theorem]{Remark}
\newtheorem{notation}[theorem]{Notation}

% ── Python listings style ──────────────────────────────────────────────
\definecolor{jugeoBlue}{HTML}{2563EB}
\definecolor{jugeoGreen}{HTML}{059669}
\definecolor{jugeoRed}{HTML}{DC2626}
\definecolor{jugeoPurple}{HTML}{7C3AED}
\definecolor{jugeoGray}{HTML}{6B7280}
\definecolor{jugeoBg}{HTML}{F8FAFC}

\lstdefinestyle{jugeo-python}{
  language=Python,
  basicstyle=\ttfamily\small,
  keywordstyle=\color{jugeoBlue}\bfseries,
  stringstyle=\color{jugeoGreen},
  commentstyle=\color{jugeoGray}\itshape,
  numberstyle=\tiny\color{jugeoGray},
  backgroundcolor=\color{jugeoBg},
  numbers=left,
  numbersep=6pt,
  frame=single,
  framerule=0.4pt,
  rulecolor=\color{jugeoGray!40},
  breaklines=true,
  breakatwhitespace=true,
  showstringspaces=false,
  tabsize=4,
  xleftmargin=1.5em,
  framexleftmargin=1.5em,
  morekeywords={dataclass,frozen,field,Enum,IntEnum,auto,Any,
                Mapping,Sequence,Optional,match,case},
  literate={→}{{$\to$}}1 {←}{{$\leftarrow$}}1
           {≤}{{$\leq$}}1 {≥}{{$\geq$}}1 {≠}{{$\neq$}}1
           {∈}{{$\in$}}1 {∀}{{$\forall$}}1 {∃}{{$\exists$}}1
           {⊕}{{$\oplus$}}1 {⊖}{{$\ominus$}}1,
  escapeinside={(*@}{@*)},
}
\lstset{style=jugeo-python}

\lstdefinestyle{jugeo-output}{
  basicstyle=\ttfamily\small,
  backgroundcolor=\color{jugeoBg},
  frame=single,
  framerule=0.4pt,
  rulecolor=\color{jugeoGray!40},
  breaklines=true,
  showstringspaces=false,
  xleftmargin=1.5em,
  framexleftmargin=0.5em,
  numbers=none,
}

% ── JuGeo-specific macros ──────────────────────────────────────────────

% Judgment 8-tuple
\newcommand{\J}{\mathcal{J}}
\newcommand{\Jud}[8]{(#1,\,#2,\,#3,\,#4,\,#5,\,#6,\,#7,\,#8)}
\newcommand{\JudShort}[1]{\J_{#1}}

% Components
\newcommand{\coord}{\mathsf{c}}            % coordinate
\newcommand{\prop}{\varphi}                 % proposition
\newcommand{\carrier}{\mathcal{A}}          % carrier type
\newcommand{\evid}{\mathcal{E}}             % evidence bundle
\newcommand{\oblig}{\mathcal{O}}            % obligations
\newcommand{\obstr}{\mathcal{B}}            % obstructions
\newcommand{\trust}{\mathcal{T}}            % trust annotation
\newcommand{\prov}{\Pi}                     % provenance

% Trust algebra
\newcommand{\Talg}{\mathfrak{T}}            % trust ordered algebra
\newcommand{\Eadm}{\mathcal{E}_{\mathrm{adm}}}
\newcommand{\tleq}{\preceq}                 % trust ordering
\newcommand{\tjoin}{\oplus}                 % conservative join
\newcommand{\tatten}{\ominus}               % attenuation
\newcommand{\tpromote}{\uparrow_\pi}        % promotion
\newcommand{\tdemote}{\downarrow_\chi}       % demotion

% Trust tiers
\newcommand{\Tcontra}{\ensuremath{\mathsf{CONTRADICTED}}}
\newcommand{\Tunver}{\ensuremath{\mathsf{UNVERIFIED}}}
\newcommand{\Tcopilot}{\ensuremath{\mathsf{COPILOT}}}
\newcommand{\Toracle}{\ensuremath{\mathsf{ORACLE}}}
\newcommand{\Truntime}{\ensuremath{\mathsf{RUNTIME}}}
\newcommand{\Tsolver}{\ensuremath{\mathsf{SOLVER}}}
\newcommand{\Tproof}{\ensuremath{\mathsf{PROOF}}}

% Site and geometry
\newcommand{\Site}{\mathcal{S}}             % semantic site
\newcommand{\Cov}{\mathrm{Cov}}            % covering family
\newcommand{\Ob}{\mathrm{Ob}}              % objects
\newcommand{\Mor}{\mathrm{Mor}}            % morphisms
\newcommand{\res}{\mathrm{res}}            % restriction
\newcommand{\Cech}{\check{C}}              % Čech complex
\newcommand{\Hcoh}[1]{H^{#1}}             % cohomology group

% Descent
\newcommand{\descent}{\mathrm{desc}}
\newcommand{\glue}{\mathrm{glue}}
\newcommand{\overlap}[2]{#1 \cap #2}

% Semantic moves
\newcommand{\VERIFY}{\textsc{Verify}}
\newcommand{\CONSTRUCT}{\textsc{Construct}}
\newcommand{\NORMALIZE}{\textsc{Normalize}}
\newcommand{\GLUE}{\textsc{Glue}}
\newcommand{\RESTRICT}{\textsc{Restrict}}
\newcommand{\TRANSPORT}{\textsc{Transport}}
\newcommand{\REPAIR}{\textsc{Repair}}
\newcommand{\PROMOTE}{\textsc{Promote}}
\newcommand{\CHALLENGE}{\textsc{Challenge}}
\newcommand{\ESCALATE}{\textsc{Escalate}}

% Fragments
\newcommand{\QFLIA}{\texttt{QF\_LIA}}
\newcommand{\QFLRA}{\texttt{QF\_LRA}}
\newcommand{\QFBV}{\texttt{QF\_BV}}
\newcommand{\QFUF}{\texttt{QF\_UF}}

% Misc
\newcommand{\Python}{\textsc{Python}}
\newcommand{\JuGeo}{\textsc{JuGeo}}
\newcommand{\LEAN}{\textsc{Lean}\,4}
\newcommand{\Fstar}{F$^{\star}$}
\newcommand{\Coq}{\textsc{Coq}}
\newcommand{\Dafny}{\textsc{Dafny}}

% Common shortcuts
\newcommand{\ie}{i.e.\@\xspace}
\newcommand{\eg}{e.g.\@\xspace}
\newcommand{\cf}{cf.\@\xspace}
\newcommand{\wrt}{w.r.t.\@\xspace}

% Evaluation shortcuts
\newcommand{\numcases}{300}
\newcommand{\numspec}{100}
\newcommand{\numeq}{100}
\newcommand{\numbug}{100}
\newcommand{\medtime}{6.5\,ms}
\newcommand{\totaltime}{1.949\,s}


% ── Packages ────────────────────────────────────────────────────────────
\usepackage{amsmath,amssymb,amsthm,mathtools}
\usepackage{tikz-cd}
\usepackage{listings}
\usepackage{xcolor}
\usepackage{xspace}
\usepackage{booktabs}
\usepackage{multirow}
\usepackage{enumitem}
\usepackage[expansion=false]{microtype}
\usepackage{hyperref}
\usepackage{cleveref}
% \usepackage{thmtools}  % disabled: not available in this TeX installation
\usepackage{float}
\usepackage{caption}
\usepackage{subcaption}
\usepackage{tabularx}
\usepackage{stmaryrd}       % \llbracket, \rrbracket
\usepackage{bbm}            % \mathbbm{1}

% ── Page geometry ───────────────────────────────────────────────────────
\usepackage[margin=1in]{geometry}

% ── Theorem environments ────────────────────────────────────────────────
\theoremstyle{plain}
\newtheorem{theorem}{Theorem}[section]
\newtheorem{lemma}[theorem]{Lemma}
\newtheorem{proposition}[theorem]{Proposition}
\newtheorem{corollary}[theorem]{Corollary}

\theoremstyle{definition}
\newtheorem{definition}[theorem]{Definition}
\newtheorem{example}[theorem]{Example}
\newtheorem{construction}[theorem]{Construction}

\theoremstyle{remark}
\newtheorem{remark}[theorem]{Remark}
\newtheorem{notation}[theorem]{Notation}

% ── Python listings style ──────────────────────────────────────────────
\definecolor{jugeoBlue}{HTML}{2563EB}
\definecolor{jugeoGreen}{HTML}{059669}
\definecolor{jugeoRed}{HTML}{DC2626}
\definecolor{jugeoPurple}{HTML}{7C3AED}
\definecolor{jugeoGray}{HTML}{6B7280}
\definecolor{jugeoBg}{HTML}{F8FAFC}

\lstdefinestyle{jugeo-python}{
  language=Python,
  basicstyle=\ttfamily\small,
  keywordstyle=\color{jugeoBlue}\bfseries,
  stringstyle=\color{jugeoGreen},
  commentstyle=\color{jugeoGray}\itshape,
  numberstyle=\tiny\color{jugeoGray},
  backgroundcolor=\color{jugeoBg},
  numbers=left,
  numbersep=6pt,
  frame=single,
  framerule=0.4pt,
  rulecolor=\color{jugeoGray!40},
  breaklines=true,
  breakatwhitespace=true,
  showstringspaces=false,
  tabsize=4,
  xleftmargin=1.5em,
  framexleftmargin=1.5em,
  morekeywords={dataclass,frozen,field,Enum,IntEnum,auto,Any,
                Mapping,Sequence,Optional,match,case},
  literate={→}{{$\to$}}1 {←}{{$\leftarrow$}}1
           {≤}{{$\leq$}}1 {≥}{{$\geq$}}1 {≠}{{$\neq$}}1
           {∈}{{$\in$}}1 {∀}{{$\forall$}}1 {∃}{{$\exists$}}1
           {⊕}{{$\oplus$}}1 {⊖}{{$\ominus$}}1,
  escapeinside={(*@}{@*)},
}
\lstset{style=jugeo-python}

\lstdefinestyle{jugeo-output}{
  basicstyle=\ttfamily\small,
  backgroundcolor=\color{jugeoBg},
  frame=single,
  framerule=0.4pt,
  rulecolor=\color{jugeoGray!40},
  breaklines=true,
  showstringspaces=false,
  xleftmargin=1.5em,
  framexleftmargin=0.5em,
  numbers=none,
}

% ── JuGeo-specific macros ──────────────────────────────────────────────

% Judgment 8-tuple
\newcommand{\J}{\mathcal{J}}
\newcommand{\Jud}[8]{(#1,\,#2,\,#3,\,#4,\,#5,\,#6,\,#7,\,#8)}
\newcommand{\JudShort}[1]{\J_{#1}}

% Components
\newcommand{\coord}{\mathsf{c}}            % coordinate
\newcommand{\prop}{\varphi}                 % proposition
\newcommand{\carrier}{\mathcal{A}}          % carrier type
\newcommand{\evid}{\mathcal{E}}             % evidence bundle
\newcommand{\oblig}{\mathcal{O}}            % obligations
\newcommand{\obstr}{\mathcal{B}}            % obstructions
\newcommand{\trust}{\mathcal{T}}            % trust annotation
\newcommand{\prov}{\Pi}                     % provenance

% Trust algebra
\newcommand{\Talg}{\mathfrak{T}}            % trust ordered algebra
\newcommand{\Eadm}{\mathcal{E}_{\mathrm{adm}}}
\newcommand{\tleq}{\preceq}                 % trust ordering
\newcommand{\tjoin}{\oplus}                 % conservative join
\newcommand{\tatten}{\ominus}               % attenuation
\newcommand{\tpromote}{\uparrow_\pi}        % promotion
\newcommand{\tdemote}{\downarrow_\chi}       % demotion

% Trust tiers
\newcommand{\Tcontra}{\ensuremath{\mathsf{CONTRADICTED}}}
\newcommand{\Tunver}{\ensuremath{\mathsf{UNVERIFIED}}}
\newcommand{\Tcopilot}{\ensuremath{\mathsf{COPILOT}}}
\newcommand{\Toracle}{\ensuremath{\mathsf{ORACLE}}}
\newcommand{\Truntime}{\ensuremath{\mathsf{RUNTIME}}}
\newcommand{\Tsolver}{\ensuremath{\mathsf{SOLVER}}}
\newcommand{\Tproof}{\ensuremath{\mathsf{PROOF}}}

% Site and geometry
\newcommand{\Site}{\mathcal{S}}             % semantic site
\newcommand{\Cov}{\mathrm{Cov}}            % covering family
\newcommand{\Ob}{\mathrm{Ob}}              % objects
\newcommand{\Mor}{\mathrm{Mor}}            % morphisms
\newcommand{\res}{\mathrm{res}}            % restriction
\newcommand{\Cech}{\check{C}}              % Čech complex
\newcommand{\Hcoh}[1]{H^{#1}}             % cohomology group

% Descent
\newcommand{\descent}{\mathrm{desc}}
\newcommand{\glue}{\mathrm{glue}}
\newcommand{\overlap}[2]{#1 \cap #2}

% Semantic moves
\newcommand{\VERIFY}{\textsc{Verify}}
\newcommand{\CONSTRUCT}{\textsc{Construct}}
\newcommand{\NORMALIZE}{\textsc{Normalize}}
\newcommand{\GLUE}{\textsc{Glue}}
\newcommand{\RESTRICT}{\textsc{Restrict}}
\newcommand{\TRANSPORT}{\textsc{Transport}}
\newcommand{\REPAIR}{\textsc{Repair}}
\newcommand{\PROMOTE}{\textsc{Promote}}
\newcommand{\CHALLENGE}{\textsc{Challenge}}
\newcommand{\ESCALATE}{\textsc{Escalate}}

% Fragments
\newcommand{\QFLIA}{\texttt{QF\_LIA}}
\newcommand{\QFLRA}{\texttt{QF\_LRA}}
\newcommand{\QFBV}{\texttt{QF\_BV}}
\newcommand{\QFUF}{\texttt{QF\_UF}}

% Misc
\newcommand{\Python}{\textsc{Python}}
\newcommand{\JuGeo}{\textsc{JuGeo}}
\newcommand{\LEAN}{\textsc{Lean}\,4}
\newcommand{\Fstar}{F$^{\star}$}
\newcommand{\Coq}{\textsc{Coq}}
\newcommand{\Dafny}{\textsc{Dafny}}

% Common shortcuts
\newcommand{\ie}{i.e.\@\xspace}
\newcommand{\eg}{e.g.\@\xspace}
\newcommand{\cf}{cf.\@\xspace}
\newcommand{\wrt}{w.r.t.\@\xspace}

% Evaluation shortcuts
\newcommand{\numcases}{300}
\newcommand{\numspec}{100}
\newcommand{\numeq}{100}
\newcommand{\numbug}{100}
\newcommand{\medtime}{6.5\,ms}
\newcommand{\totaltime}{1.949\,s}


% ── Packages ────────────────────────────────────────────────────────────
\usepackage{amsmath,amssymb,amsthm,mathtools}
\usepackage{tikz-cd}
\usepackage{listings}
\usepackage{xcolor}
\usepackage{xspace}
\usepackage{booktabs}
\usepackage{multirow}
\usepackage{enumitem}
\usepackage[expansion=false]{microtype}
\usepackage{hyperref}
\usepackage{cleveref}
% \usepackage{thmtools}  % disabled: not available in this TeX installation
\usepackage{float}
\usepackage{caption}
\usepackage{subcaption}
\usepackage{tabularx}
\usepackage{stmaryrd}       % \llbracket, \rrbracket
\usepackage{bbm}            % \mathbbm{1}

% ── Page geometry ───────────────────────────────────────────────────────
\usepackage[margin=1in]{geometry}

% ── Theorem environments ────────────────────────────────────────────────
\theoremstyle{plain}
\newtheorem{theorem}{Theorem}[section]
\newtheorem{lemma}[theorem]{Lemma}
\newtheorem{proposition}[theorem]{Proposition}
\newtheorem{corollary}[theorem]{Corollary}

\theoremstyle{definition}
\newtheorem{definition}[theorem]{Definition}
\newtheorem{example}[theorem]{Example}
\newtheorem{construction}[theorem]{Construction}

\theoremstyle{remark}
\newtheorem{remark}[theorem]{Remark}
\newtheorem{notation}[theorem]{Notation}

% ── Python listings style ──────────────────────────────────────────────
\definecolor{jugeoBlue}{HTML}{2563EB}
\definecolor{jugeoGreen}{HTML}{059669}
\definecolor{jugeoRed}{HTML}{DC2626}
\definecolor{jugeoPurple}{HTML}{7C3AED}
\definecolor{jugeoGray}{HTML}{6B7280}
\definecolor{jugeoBg}{HTML}{F8FAFC}

\lstdefinestyle{jugeo-python}{
  language=Python,
  basicstyle=\ttfamily\small,
  keywordstyle=\color{jugeoBlue}\bfseries,
  stringstyle=\color{jugeoGreen},
  commentstyle=\color{jugeoGray}\itshape,
  numberstyle=\tiny\color{jugeoGray},
  backgroundcolor=\color{jugeoBg},
  numbers=left,
  numbersep=6pt,
  frame=single,
  framerule=0.4pt,
  rulecolor=\color{jugeoGray!40},
  breaklines=true,
  breakatwhitespace=true,
  showstringspaces=false,
  tabsize=4,
  xleftmargin=1.5em,
  framexleftmargin=1.5em,
  morekeywords={dataclass,frozen,field,Enum,IntEnum,auto,Any,
                Mapping,Sequence,Optional,match,case},
  literate={→}{{$\to$}}1 {←}{{$\leftarrow$}}1
           {≤}{{$\leq$}}1 {≥}{{$\geq$}}1 {≠}{{$\neq$}}1
           {∈}{{$\in$}}1 {∀}{{$\forall$}}1 {∃}{{$\exists$}}1
           {⊕}{{$\oplus$}}1 {⊖}{{$\ominus$}}1,
  escapeinside={(*@}{@*)},
}
\lstset{style=jugeo-python}

\lstdefinestyle{jugeo-output}{
  basicstyle=\ttfamily\small,
  backgroundcolor=\color{jugeoBg},
  frame=single,
  framerule=0.4pt,
  rulecolor=\color{jugeoGray!40},
  breaklines=true,
  showstringspaces=false,
  xleftmargin=1.5em,
  framexleftmargin=0.5em,
  numbers=none,
}

% ── JuGeo-specific macros ──────────────────────────────────────────────

% Judgment 8-tuple
\newcommand{\J}{\mathcal{J}}
\newcommand{\Jud}[8]{(#1,\,#2,\,#3,\,#4,\,#5,\,#6,\,#7,\,#8)}
\newcommand{\JudShort}[1]{\J_{#1}}

% Components
\newcommand{\coord}{\mathsf{c}}            % coordinate
\newcommand{\prop}{\varphi}                 % proposition
\newcommand{\carrier}{\mathcal{A}}          % carrier type
\newcommand{\evid}{\mathcal{E}}             % evidence bundle
\newcommand{\oblig}{\mathcal{O}}            % obligations
\newcommand{\obstr}{\mathcal{B}}            % obstructions
\newcommand{\trust}{\mathcal{T}}            % trust annotation
\newcommand{\prov}{\Pi}                     % provenance

% Trust algebra
\newcommand{\Talg}{\mathfrak{T}}            % trust ordered algebra
\newcommand{\Eadm}{\mathcal{E}_{\mathrm{adm}}}
\newcommand{\tleq}{\preceq}                 % trust ordering
\newcommand{\tjoin}{\oplus}                 % conservative join
\newcommand{\tatten}{\ominus}               % attenuation
\newcommand{\tpromote}{\uparrow_\pi}        % promotion
\newcommand{\tdemote}{\downarrow_\chi}       % demotion

% Trust tiers
\newcommand{\Tcontra}{\ensuremath{\mathsf{CONTRADICTED}}}
\newcommand{\Tunver}{\ensuremath{\mathsf{UNVERIFIED}}}
\newcommand{\Tcopilot}{\ensuremath{\mathsf{COPILOT}}}
\newcommand{\Toracle}{\ensuremath{\mathsf{ORACLE}}}
\newcommand{\Truntime}{\ensuremath{\mathsf{RUNTIME}}}
\newcommand{\Tsolver}{\ensuremath{\mathsf{SOLVER}}}
\newcommand{\Tproof}{\ensuremath{\mathsf{PROOF}}}

% Site and geometry
\newcommand{\Site}{\mathcal{S}}             % semantic site
\newcommand{\Cov}{\mathrm{Cov}}            % covering family
\newcommand{\Ob}{\mathrm{Ob}}              % objects
\newcommand{\Mor}{\mathrm{Mor}}            % morphisms
\newcommand{\res}{\mathrm{res}}            % restriction
\newcommand{\Cech}{\check{C}}              % Čech complex
\newcommand{\Hcoh}[1]{H^{#1}}             % cohomology group

% Descent
\newcommand{\descent}{\mathrm{desc}}
\newcommand{\glue}{\mathrm{glue}}
\newcommand{\overlap}[2]{#1 \cap #2}

% Semantic moves
\newcommand{\VERIFY}{\textsc{Verify}}
\newcommand{\CONSTRUCT}{\textsc{Construct}}
\newcommand{\NORMALIZE}{\textsc{Normalize}}
\newcommand{\GLUE}{\textsc{Glue}}
\newcommand{\RESTRICT}{\textsc{Restrict}}
\newcommand{\TRANSPORT}{\textsc{Transport}}
\newcommand{\REPAIR}{\textsc{Repair}}
\newcommand{\PROMOTE}{\textsc{Promote}}
\newcommand{\CHALLENGE}{\textsc{Challenge}}
\newcommand{\ESCALATE}{\textsc{Escalate}}

% Fragments
\newcommand{\QFLIA}{\texttt{QF\_LIA}}
\newcommand{\QFLRA}{\texttt{QF\_LRA}}
\newcommand{\QFBV}{\texttt{QF\_BV}}
\newcommand{\QFUF}{\texttt{QF\_UF}}

% Misc
\newcommand{\Python}{\textsc{Python}}
\newcommand{\JuGeo}{\textsc{JuGeo}}
\newcommand{\LEAN}{\textsc{Lean}\,4}
\newcommand{\Fstar}{F$^{\star}$}
\newcommand{\Coq}{\textsc{Coq}}
\newcommand{\Dafny}{\textsc{Dafny}}

% Common shortcuts
\newcommand{\ie}{i.e.\@\xspace}
\newcommand{\eg}{e.g.\@\xspace}
\newcommand{\cf}{cf.\@\xspace}
\newcommand{\wrt}{w.r.t.\@\xspace}

% Evaluation shortcuts
\newcommand{\numcases}{300}
\newcommand{\numspec}{100}
\newcommand{\numeq}{100}
\newcommand{\numbug}{100}
\newcommand{\medtime}{6.5\,ms}
\newcommand{\totaltime}{1.949\,s}

% experiment-data.tex — AUTO-GENERATED from real jugeo CLI runs
% DO NOT EDIT — regenerate with: python3 experiments/run_universal_metrics.py
% Generated from 30 programs across 6 domains

\newcommand{\expTotalPrograms}{30}
\newcommand{\expVerified}{30}
\newcommand{\expAccuracy}{100.0\%}
\newcommand{\expCoordsMin}{2}
\newcommand{\expCoordsMax}{10}
\newcommand{\expCoordsMean}{5.2}
\newcommand{\expCoordsMedian}{5.5}
\newcommand{\expPropsMin}{4}
\newcommand{\expPropsMax}{20}
\newcommand{\expPropsMean}{10.4}
\newcommand{\expPropsSum}{311}
\newcommand{\expPropsOkSum}{310}
\newcommand{\expTimeMin}{3.506\,s}
\newcommand{\expTimeMax}{6.714\,s}
\newcommand{\expTimeMean}{4.64\,s}
\newcommand{\expTimeMedian}{4.508\,s}
\newcommand{\expTimeTotal}{139.204\,s}
\newcommand{\expObstructionTotal}{0}
\newcommand{\expHOneZero}{30}
\newcommand{\expDomainCount}{6}
\newcommand{\expTrustCOPILOTSUGGESTED}{30}
\newcommand{\expDomainDsCount}{5}
\newcommand{\expDomainDsVerified}{5}
\newcommand{\expDomainDsAvgCoords}{7.6}
\newcommand{\expDomainDsAvgProps}{15.2}
\newcommand{\expDomainMathCount}{5}
\newcommand{\expDomainMathVerified}{5}
\newcommand{\expDomainMathAvgCoords}{4.8}
\newcommand{\expDomainMathAvgProps}{9.4}
\newcommand{\expDomainSmCount}{5}
\newcommand{\expDomainSmVerified}{5}
\newcommand{\expDomainSmAvgCoords}{7.6}
\newcommand{\expDomainSmAvgProps}{15.2}
\newcommand{\expDomainSortCount}{5}
\newcommand{\expDomainSortVerified}{5}
\newcommand{\expDomainSortAvgCoords}{2.4}
\newcommand{\expDomainSortAvgProps}{4.8}
\newcommand{\expDomainStrCount}{5}
\newcommand{\expDomainStrVerified}{5}
\newcommand{\expDomainStrAvgCoords}{3.2}
\newcommand{\expDomainStrAvgProps}{6.4}
\newcommand{\expDomainWebCount}{5}
\newcommand{\expDomainWebVerified}{5}
\newcommand{\expDomainWebAvgCoords}{5.6}
\newcommand{\expDomainWebAvgProps}{11.2}

% subsystem-data.tex — AUTO-GENERATED from real jugeo subsystem experiments
% DO NOT EDIT — regenerate with: python3 experiments/run_subsystem_experiments.py

\newcommand{\subCliTotal}{10}
\newcommand{\subCliVerified}{9}
\newcommand{\subCliAccuracy}{90.0\%}
\newcommand{\subCliMeanTime}{4.132\,s}
\newcommand{\subCliMinTime}{3.472\,s}
\newcommand{\subCliMaxTime}{5.372\,s}
\newcommand{\subCliTotalCoords}{54}
\newcommand{\subCliTotalProps}{107}
\newcommand{\subCliTotalPropsOk}{106}
\newcommand{\subSiteCalls}{90}
\newcommand{\subSiteSuccessful}{69}
\newcommand{\subSiteSuccessRate}{76.7\%}
\newcommand{\subSiteMeanTime}{0.0001\,s}
\newcommand{\subSiteTotalTime}{0.005\,s}
\newcommand{\subSpecificationsatisfactionSmallTime}{0.0003\,s}
\newcommand{\subSpecificationsatisfactionMediumTime}{0.0001\,s}
\newcommand{\subSpecificationsatisfactionLargeTime}{0.0001\,s}
\newcommand{\subInhabitantfleetSmallTime}{0.0\,s}
\newcommand{\subInhabitantfleetMediumTime}{0.0\,s}
\newcommand{\subInhabitantfleetLargeTime}{0.0\,s}
\newcommand{\subTheoremecologySmallTime}{0.0\,s}
\newcommand{\subTheoremecologyMediumTime}{0.0\,s}
\newcommand{\subTheoremecologyLargeTime}{0.0\,s}
\newcommand{\subStatespaceexplorationSmallTime}{0.0\,s}
\newcommand{\subStatespaceexplorationMediumTime}{0.0\,s}
\newcommand{\subStatespaceexplorationLargeTime}{0.0\,s}
\newcommand{\subRepairsemanticsSmallTime}{0.0\,s}
\newcommand{\subRepairsemanticsMediumTime}{0.0\,s}
\newcommand{\subRepairsemanticsLargeTime}{0.0\,s}
\newcommand{\subReplaygluingSmallTime}{0.0\,s}
\newcommand{\subReplaygluingMediumTime}{0.0\,s}
\newcommand{\subReplaygluingLargeTime}{0.0\,s}
\newcommand{\subSemanticfuturesSmallTime}{0.0\,s}
\newcommand{\subSemanticfuturesMediumTime}{0.0\,s}
\newcommand{\subSemanticfuturesLargeTime}{0.0\,s}
\newcommand{\subHypercovertreatySmallTime}{0.0\,s}
\newcommand{\subHypercovertreatyMediumTime}{0.0\,s}
\newcommand{\subHypercovertreatyLargeTime}{0.0\,s}
\newcommand{\subEncodeforsolverSmallTime}{0.0026\,s}
\newcommand{\subEncodeforsolverMediumTime}{0.0002\,s}
\newcommand{\subEncodeforsolverLargeTime}{0.0001\,s}
\newcommand{\subInterfaceroutingSmallTime}{0.0\,s}
\newcommand{\subInterfaceroutingMediumTime}{0.0\,s}
\newcommand{\subInterfaceroutingLargeTime}{0.0\,s}
\newcommand{\subOrchestrateverificationSmallTime}{0.0002\,s}
\newcommand{\subOrchestrateverificationMediumTime}{0.0\,s}
\newcommand{\subOrchestrateverificationLargeTime}{0.0\,s}
\newcommand{\subRunfulldescentSmallTime}{0.0\,s}
\newcommand{\subRunfulldescentMediumTime}{0.0\,s}
\newcommand{\subRunfulldescentLargeTime}{0.0\,s}
\newcommand{\subKernellifecycleSmallTime}{0.0\,s}
\newcommand{\subKernellifecycleMediumTime}{0.0\,s}
\newcommand{\subKernellifecycleLargeTime}{0.0\,s}
\newcommand{\subDiscoverypipelineSmallTime}{0.0\,s}
\newcommand{\subDiscoverypipelineMediumTime}{0.0\,s}
\newcommand{\subDiscoverypipelineLargeTime}{0.0\,s}
\newcommand{\subAnalogytransportSmallTime}{0.0\,s}
\newcommand{\subAnalogytransportMediumTime}{0.0\,s}
\newcommand{\subAnalogytransportLargeTime}{0.0\,s}
\newcommand{\subFormalcoresiteSmallTime}{0.0001\,s}
\newcommand{\subFormalcoresiteMediumTime}{0.0\,s}
\newcommand{\subFormalcoresiteLargeTime}{0.0\,s}
\newcommand{\subProblematlasSmallTime}{0.0\,s}
\newcommand{\subProblematlasMediumTime}{0.0\,s}
\newcommand{\subProblematlasLargeTime}{0.0\,s}
\newcommand{\subPublicalignmentSmallTime}{0.0\,s}
\newcommand{\subPublicalignmentMediumTime}{0.0\,s}
\newcommand{\subPublicalignmentLargeTime}{0.0\,s}
\newcommand{\subRegimebootstrappingSmallTime}{0.0\,s}
\newcommand{\subRegimebootstrappingMediumTime}{0.0\,s}
\newcommand{\subRegimebootstrappingLargeTime}{0.0\,s}
\newcommand{\subRelationalrefinementSmallTime}{0.0\,s}
\newcommand{\subRelationalrefinementMediumTime}{0.0\,s}
\newcommand{\subRelationalrefinementLargeTime}{0.0\,s}
\newcommand{\subSemanticclosureSmallTime}{0.0\,s}
\newcommand{\subSemanticclosureMediumTime}{0.0\,s}
\newcommand{\subSemanticclosureLargeTime}{0.0\,s}
\newcommand{\subTheoremeconomicsSmallTime}{0.0001\,s}
\newcommand{\subTheoremeconomicsMediumTime}{0.0\,s}
\newcommand{\subTheoremeconomicsLargeTime}{0.0\,s}
\newcommand{\subBenchmarksuiteSmallTime}{0.0\,s}
\newcommand{\subBenchmarksuiteMediumTime}{0.0\,s}
\newcommand{\subBenchmarksuiteLargeTime}{0.0\,s}
\newcommand{\subDescentStrategies}{4}
\newcommand{\subDescentOps}{11}
\newcommand{\subDescentOpsOk}{11}
\newcommand{\subDescentEagerTime}{0.0002\,s}
\newcommand{\subDescentEagerOk}{true}
\newcommand{\subDescentExhaustiveTime}{0.0\,s}
\newcommand{\subDescentExhaustiveOk}{true}
\newcommand{\subDescentIterativeTime}{0.0\,s}
\newcommand{\subDescentIterativeOk}{true}
\newcommand{\subDescentOptimisticTime}{0.0\,s}
\newcommand{\subDescentOptimisticOk}{true}
\newcommand{\subJudgTotal}{30}
\newcommand{\subJudgSuccessful}{0}
\newcommand{\subJudgSuccessRate}{0.0\%}
\newcommand{\subJudgMeanTimeUs}{7.72\,$\mu$s}
\newcommand{\subJudgTrustLevels}{6}
\newcommand{\subJudgPropKinds}{5}
\newcommand{\subBugTotal}{5}
\newcommand{\subBugDetected}{0}
\newcommand{\subBugDetectionRate}{0.0\%}
\newcommand{\subBugFalsePositiveRate}{0.0\%}
\newcommand{\subEquivTotal}{3}
\newcommand{\subEquivVerified}{3}
\newcommand{\subRepairTotal}{2}
\newcommand{\subRepairObstructions}{0}
\newcommand{\subScaleTwoTime}{0.0001\,s}
\newcommand{\subScaleFourTime}{0.0\,s}
\newcommand{\subScaleEightTime}{0.0\,s}
\newcommand{\subScaleTwelveTime}{0.0\,s}
\newcommand{\subScaleSixteenTime}{0.0\,s}
\newcommand{\subScaleTwentyTime}{0.0\,s}
\newcommand{\subTotalExperimentTime}{95.6\,s}

% comprehensive-data.tex — AUTO-GENERATED from real jugeo experiments
% DO NOT EDIT — regenerate with: python3 experiments/run_comprehensive_experiments.py
% Generated: 2026-03-23 09:19:06

% --- Size scaling metrics ---
\newcommand{\compTotalPrograms}{18}
\newcommand{\compTotalVerified}{18}
\newcommand{\compOverallAccuracy}{100.0\%}
\newcommand{\compTinyCount}{5}
\newcommand{\compTinyVerified}{5}
\newcommand{\compTinyMeanCoords}{3}
\newcommand{\compTinyMeanProps}{6}
\newcommand{\compTinyMeanTime}{4.95\,s}
\newcommand{\compTinyMeanLines}{5}
\newcommand{\compTinyMinTime}{4.45\,s}
\newcommand{\compTinyMaxTime}{5.51\,s}
\newcommand{\compSmallCount}{5}
\newcommand{\compSmallVerified}{5}
\newcommand{\compSmallMeanCoords}{3.6}
\newcommand{\compSmallMeanProps}{7.2}
\newcommand{\compSmallMeanTime}{4.85\,s}
\newcommand{\compSmallMeanLines}{16.0}
\newcommand{\compSmallMinTime}{4.30\,s}
\newcommand{\compSmallMaxTime}{5.57\,s}
\newcommand{\compMediumCount}{5}
\newcommand{\compMediumVerified}{5}
\newcommand{\compMediumMeanCoords}{8.6}
\newcommand{\compMediumMeanProps}{17.2}
\newcommand{\compMediumMeanTime}{4.47\,s}
\newcommand{\compMediumMeanLines}{45.0}
\newcommand{\compMediumMinTime}{4.26\,s}
\newcommand{\compMediumMaxTime}{4.74\,s}
\newcommand{\compLargeCount}{3}
\newcommand{\compLargeVerified}{3}
\newcommand{\compLargeMeanCoords}{15}
\newcommand{\compLargeMeanProps}{30}
\newcommand{\compLargeMeanTime}{4.31\,s}
\newcommand{\compLargeMeanLines}{79.0}
\newcommand{\compLargeMinTime}{4.27\,s}
\newcommand{\compLargeMaxTime}{4.38\,s}
% --- Aggregate timing ---
\newcommand{\compTimeMean}{4.68\,s}
\newcommand{\compTimeMedian}{4.50\,s}
\newcommand{\compTimeMin}{4.265\,s}
\newcommand{\compTimeMax}{5.571\,s}
\newcommand{\compTimeTotal}{84.3\,s}
\newcommand{\compTimeStdev}{0.45\,s}
% --- Aggregate structure ---
\newcommand{\compCoordsMean}{6.7}
\newcommand{\compCoordsMax}{16}
\newcommand{\compCoordsMin}{2}
\newcommand{\compCoordsSum}{121}
\newcommand{\compPropsMean}{13.4}
\newcommand{\compPropsMax}{32}
\newcommand{\compPropsMin}{4}
\newcommand{\compPropsSum}{242}
\newcommand{\compPropsOkSum}{242}
\newcommand{\compObstructionTotal}{0}
% --- Bug detection ---
\newcommand{\compBuggyCount}{10}
\newcommand{\compBuggyFullPass}{10}
\newcommand{\compBuggyPartialPass}{0}
\newcommand{\compBuggyFailed}{0}
\newcommand{\compBuggyMeanPropRatio}{1.00}
\newcommand{\compCorrectMeanPropRatio}{1.00}
\newcommand{\compBuggyMeanTime}{5.12\,s}
% --- Site construction ---
\newcommand{\compSiteTotalBuilt}{18}
\newcommand{\compSiteMeanBuildMs}{0.05}
\newcommand{\compSiteMaxBuildMs}{0.14}
\newcommand{\compSiteMeanCoords}{6.5}
\newcommand{\compSiteMeanMorphisms}{14.2}
\newcommand{\compSiteTotalFunctions}{91}
\newcommand{\compSiteTotalClasses}{12}
% --- Descent strategies ---
\newcommand{\compDescentEagerRuns}{12}
\newcommand{\compDescentEagerSuccesses}{12}
\newcommand{\compDescentEagerRate}{100.0\%}
\newcommand{\compDescentEagerMeanMs}{0.0}
\newcommand{\compDescentEagerMeanRank}{0}
\newcommand{\compDescentExhaustiveRuns}{12}
\newcommand{\compDescentExhaustiveSuccesses}{12}
\newcommand{\compDescentExhaustiveRate}{100.0\%}
\newcommand{\compDescentExhaustiveMeanMs}{0.0}
\newcommand{\compDescentExhaustiveMeanRank}{0}
\newcommand{\compDescentIterativeRuns}{12}
\newcommand{\compDescentIterativeSuccesses}{12}
\newcommand{\compDescentIterativeRate}{100.0\%}
\newcommand{\compDescentIterativeMeanMs}{0.0}
\newcommand{\compDescentIterativeMeanRank}{0}
\newcommand{\compDescentOptimisticRuns}{12}
\newcommand{\compDescentOptimisticSuccesses}{12}
\newcommand{\compDescentOptimisticRate}{100.0\%}
\newcommand{\compDescentOptimisticMeanMs}{0.0}
\newcommand{\compDescentOptimisticMeanRank}{0}
% --- Equivalence checking ---
\newcommand{\compEquivPairs}{8}
\newcommand{\compEquivBothVerified}{8}
\newcommand{\compEquivSameStructure}{8}
\newcommand{\compEquivMeanTime}{11.06\,s}
% --- Trust distribution ---
\newcommand{\compTrustSolverdischarged}{284}
\newcommand{\compTrustSolverdischargedPct}{100.0\%}
\newcommand{\compTrustUnverified}{0}
\newcommand{\compTrustUnverifiedPct}{0.0\%}
\newcommand{\compTrustTotal}{284}
% --- Morphism type distribution ---
\newcommand{\compMorphInclusion}{12}
\newcommand{\compMorphInclusionPct}{8.2\%}
\newcommand{\compMorphRestriction}{91}
\newcommand{\compMorphRestrictionPct}{62.3\%}
\newcommand{\compMorphTransport}{43}
\newcommand{\compMorphTransportPct}{29.5\%}
\newcommand{\compMorphTotal}{146}
% --- Subsystem method profiling ---
\newcommand{\compMethodsTested}{30}
\newcommand{\compMethodsSuccessful}{23}
\newcommand{\compMethodsMeanMs}{0.02}
\newcommand{\compProfAnalogyTransportMeanMs}{0.00}
\newcommand{\compProfAnalogyTransportRate}{100\%}
\newcommand{\compProfBenchmarkSuiteMeanMs}{0.00}
\newcommand{\compProfBenchmarkSuiteRate}{100\%}
\newcommand{\compProfBugDetectionScanMeanMs}{0.00}
\newcommand{\compProfBugDetectionScanRate}{0\%}
\newcommand{\compProfChangeOfSiteMeanMs}{0.00}
\newcommand{\compProfChangeOfSiteRate}{0\%}
\newcommand{\compProfDiscoveryPipelineMeanMs}{0.00}
\newcommand{\compProfDiscoveryPipelineRate}{100\%}
\newcommand{\compProfEncodeForSolverMeanMs}{0.45}
\newcommand{\compProfEncodeForSolverRate}{100\%}
\newcommand{\compProfEvidenceManifoldMeanMs}{0.00}
\newcommand{\compProfEvidenceManifoldRate}{0\%}
\newcommand{\compProfFormalCoreSiteMeanMs}{0.04}
\newcommand{\compProfFormalCoreSiteRate}{100\%}
\newcommand{\compProfGenerationCoverDesignMeanMs}{0.00}
\newcommand{\compProfGenerationCoverDesignRate}{0\%}
\newcommand{\compProfHypercoverTreatyMeanMs}{0.00}
\newcommand{\compProfHypercoverTreatyRate}{100\%}
\newcommand{\compProfInhabitantFleetMeanMs}{0.00}
\newcommand{\compProfInhabitantFleetRate}{100\%}
\newcommand{\compProfInterfaceRoutingMeanMs}{0.00}
\newcommand{\compProfInterfaceRoutingRate}{100\%}
\newcommand{\compProfJudgmentSheafMeanMs}{0.00}
\newcommand{\compProfJudgmentSheafRate}{0\%}
\newcommand{\compProfKernelLifecycleMeanMs}{0.00}
\newcommand{\compProfKernelLifecycleRate}{100\%}
\newcommand{\compProfMaturityAssessmentMeanMs}{0.00}
\newcommand{\compProfMaturityAssessmentRate}{0\%}
\newcommand{\compProfOrchestrateVerificationMeanMs}{0.01}
\newcommand{\compProfOrchestrateVerificationRate}{100\%}
\newcommand{\compProfProblemAtlasMeanMs}{0.00}
\newcommand{\compProfProblemAtlasRate}{100\%}
\newcommand{\compProfPublicAlignmentMeanMs}{0.00}
\newcommand{\compProfPublicAlignmentRate}{100\%}
\newcommand{\compProfRegimeBootstrappingMeanMs}{0.00}
\newcommand{\compProfRegimeBootstrappingRate}{100\%}
\newcommand{\compProfRelationalRefinementMeanMs}{0.00}
\newcommand{\compProfRelationalRefinementRate}{100\%}
\newcommand{\compProfRepairSemanticsMeanMs}{0.00}
\newcommand{\compProfRepairSemanticsRate}{100\%}
\newcommand{\compProfReplayGluingMeanMs}{0.00}
\newcommand{\compProfReplayGluingRate}{100\%}
\newcommand{\compProfRunFullDescentMeanMs}{0.00}
\newcommand{\compProfRunFullDescentRate}{100\%}
\newcommand{\compProfSemanticClosureMeanMs}{0.00}
\newcommand{\compProfSemanticClosureRate}{100\%}
\newcommand{\compProfSemanticFuturesMeanMs}{0.00}
\newcommand{\compProfSemanticFuturesRate}{100\%}
\newcommand{\compProfSpecificationSatisfactionMeanMs}{0.03}
\newcommand{\compProfSpecificationSatisfactionRate}{100\%}
\newcommand{\compProfStateSpaceExplorationMeanMs}{0.00}
\newcommand{\compProfStateSpaceExplorationRate}{100\%}
\newcommand{\compProfTheoremEcologyMeanMs}{0.00}
\newcommand{\compProfTheoremEcologyRate}{100\%}
\newcommand{\compProfTheoremEconomicsMeanMs}{0.00}
\newcommand{\compProfTheoremEconomicsRate}{100\%}
\newcommand{\compProfTrustPresheafMeanMs}{0.00}
\newcommand{\compProfTrustPresheafRate}{0\%}

% supplementary-data.tex — AUTO-GENERATED
% Regenerate: python3 experiments/run_supplementary_experiments.py
% Generated: 2026-03-23 09:57:40

\newcommand{\suppDomainSortAvgTime}{3.59\,s}
\newcommand{\suppDomainSortMinTime}{3.18\,s}
\newcommand{\suppDomainSortMaxTime}{4.00\,s}
\newcommand{\suppDomainDsAvgTime}{3.41\,s}
\newcommand{\suppDomainDsMinTime}{3.27\,s}
\newcommand{\suppDomainDsMaxTime}{3.75\,s}
\newcommand{\suppDomainMathAvgTime}{3.76\,s}
\newcommand{\suppDomainMathMinTime}{3.15\,s}
\newcommand{\suppDomainMathMaxTime}{4.05\,s}
\newcommand{\suppDomainStrAvgTime}{3.27\,s}
\newcommand{\suppDomainStrMinTime}{3.05\,s}
\newcommand{\suppDomainStrMaxTime}{3.47\,s}
\newcommand{\suppDomainWebAvgTime}{3.33\,s}
\newcommand{\suppDomainWebMinTime}{3.25\,s}
\newcommand{\suppDomainWebMaxTime}{3.47\,s}
\newcommand{\suppDomainSmAvgTime}{3.81\,s}
\newcommand{\suppDomainSmMinTime}{3.41\,s}
\newcommand{\suppDomainSmMaxTime}{4.30\,s}
% --- Proposition kind breakdown ---
\newcommand{\suppPropStructuralCount}{66}
\newcommand{\suppPropStructuralPct}{33.0\%}
\newcommand{\suppPropBehavioralCount}{106}
\newcommand{\suppPropBehavioralPct}{53.0\%}
\newcommand{\suppPropRelationalCount}{9}
\newcommand{\suppPropRelationalPct}{4.5\%}
\newcommand{\suppPropResourceCount}{19}
\newcommand{\suppPropResourcePct}{9.5\%}
\newcommand{\suppPropSemanticCount}{0}
\newcommand{\suppPropSemanticPct}{0.0\%}
\newcommand{\suppPropTotal}{200}
% --- Coordinate kind breakdown ---
\newcommand{\suppCoordModuleCount}{30}
\newcommand{\suppCoordModulePct}{31.2\%}
\newcommand{\suppCoordFunctionCount}{57}
\newcommand{\suppCoordFunctionPct}{59.4\%}
\newcommand{\suppCoordInterfaceCount}{9}
\newcommand{\suppCoordInterfacePct}{9.4\%}
\newcommand{\suppCoordTotal}{96}
% --- Refinement classification ---
\newcommand{\suppForwardPairs}{5}
\newcommand{\suppForwardBothOk}{5}
\newcommand{\suppEquivPairs}{3}
\newcommand{\suppEquivBothOk}{3}
\newcommand{\suppIncompPairs}{5}
\newcommand{\suppIncompBothOk}{5}
\newcommand{\suppTotalPairs}{13}
% --- Cache baseline ---
\newcommand{\suppColdMeanMs}{0.663}
\newcommand{\suppWarmMeanMs}{0.019}
\newcommand{\suppCacheSpeedup}{35.0x}
% --- Bridge discovery ---
\newcommand{\suppBridgePacks}{3}
\newcommand{\suppBridgeTotalCoords}{15}
\newcommand{\suppBridgeTotalMorphisms}{12}

% \input{data-paper96}  % No data file for paper 96

\usepackage{array}
\usepackage{tikz}
\usetikzlibrary{arrows.meta,positioning,shapes.geometric,fit,calc}
\usepackage{algorithm}
\usepackage{algorithmic}

% Paper-specific macros
\newcommand{\CovAdapt}{\textsc{CopilotCoverDesignAdapter}}
\newcommand{\CovPart}{\textsc{CopilotCoverDesignParticipant}}
\newcommand{\ConstrPart}{\textsc{CopilotConstructionParticipant}}
\newcommand{\CopProp}{\textsc{CopilotProposal}}
\newcommand{\CopNeg}{\textsc{CopilotNegotiationRecord}}
\newcommand{\CopStrat}{\textsc{CopilotStrategyState}}
\newcommand{\StratAdapt}{\textsc{StrategyAdaptation}}
\newcommand{\patchcand}{\mathrm{patch\_cand}}
\newcommand{\seedcand}{\mathrm{seed\_cand}}
\newcommand{\covfam}{\mathcal{U}}
\newcommand{\glue}{\mathrm{glue}}
\newcommand{\strat}{\mathrm{strat}}

\title{LLM-Guided Cover Design:\\
Copilot Participation in Program Decomposition}
\author{JuGeo Development Team}
\date{\today}

\begin{document}
\maketitle

\begin{abstract}
JuGeo does ship a concrete cover-design integration layer in
\texttt{jugeo.generation.cover\_design.integration}, including
\CovAdapt{}.  This audited version keeps the paper focused on that
implemented adapter surface and removes unsupported claims about
benchmark-scale acceptance rates, coverage improvements, and Lean
formalization.  Where the paper discusses a broader copilot participant
pipeline, those claims are now limited to what can be directly traced to
the checked-in Python modules.
\end{abstract}

% ------------------------------------------------------------------
\section{Introduction}\label{sec:intro}
% ------------------------------------------------------------------

\paragraph{Problem statement.}
In the judgment geometry of \JuGeo{}, a program is modelled as a semantic
site $\Site$ whose objects are coordinates (functions, modules, classes) and
whose morphisms are restriction maps.  Verification proceeds by constructing
\emph{local sections}---verified judgments on individual coordinates---and
\emph{gluing} them along overlaps.  The choice of covering family
$\covfam = \{U_i\}_{i \in I}$ determines which local sections are needed
and how they compose.

A poor covering family leads to three problems: (1) uncovered coordinates
where no local section exists, (2) excessive redundancy where many sections
overlap unnecessarily, and (3) gluing failures where overlapping sections
cannot be reconciled.  Designing an optimal covering family is NP-hard in
general, motivating the use of LLM heuristics.

\paragraph{Our approach.}
We delegate the combinatorial search to GitHub Copilot, which proposes
\emph{patch candidates}---potential covering patches---that are then validated
against the site axioms.  The \CovAdapt{} class translates Copilot proposals
into the cover design sub-system's internal format, while the \CovPart{}
class manages the Copilot's participation in the design pipeline.  All
proposals are constrained to $\Tcopilot$ trust, ensuring that Copilot-designed
covers receive the same trust treatment as any other AI-generated artefact.

\paragraph{Contributions.}
\begin{enumerate}[leftmargin=2em]
  \item The \CovAdapt{} class and its three strategies: greedy, adaptive, and
        exhaustive (Definition~\ref{def:adapter}).
  \item The \CovPart{} class for managing Copilot participation in the cover
        design pipeline (§\ref{sec:participant}).
  \item Integration with \ConstrPart{} for local construction within
        Copilot-proposed patches (§\ref{sec:construction}).
  \item A failure diagnosis mechanism for rejected patches
        (Algorithm~\ref{alg:diagnosis}).
  \item An evaluation section reframed as methodology rather than as a
        source of reproduced benchmark claims.
\end{enumerate}

% ------------------------------------------------------------------
\section{Background}\label{sec:bg}
% ------------------------------------------------------------------

\subsection{Covering Families in Sheaf Theory}

\begin{definition}[Covering Family]\label{def:covering}
A \emph{covering family} for an object $U \in \Ob(\Site)$ is a collection
$\covfam = \{f_i: U_i \to U\}_{i \in I}$ of morphisms such that the
induced map $\coprod_i U_i \to U$ is surjective on the underlying topology.
The Grothendieck topology $\Cov$ specifies which collections qualify as covers.
\end{definition}

\begin{definition}[Gluing Condition]\label{def:gluing}
A presheaf $\mathcal{F}$ satisfies the \emph{gluing condition} for a cover
$\covfam = \{U_i\}$ if: given sections $s_i \in \mathcal{F}(U_i)$ that agree
on overlaps ($\res_{U_i \cap U_j}(s_i) = \res_{U_i \cap U_j}(s_j)$), there
exists a unique section $s \in \mathcal{F}(U)$ with $\res_{U_i}(s) = s_i$.
\end{definition}

\begin{definition}[\v{C}ech Cohomology]\label{def:cech}
The \v{C}ech complex $\Cech^\bullet(\covfam, \mathcal{F})$ measures the
obstruction to gluing.  The $n$-th cohomology group $\Hcoh{n}(\covfam,
\mathcal{F})$ detects $n$-fold overlap inconsistencies.  When
$\Hcoh{1}(\covfam, \mathcal{F}) = 0$, all compatible local sections can be
glued to a global section.
\end{definition}

\subsection{Patch Candidates}

A \emph{patch candidate} is a proposed element of a covering family.  It
specifies a subset of coordinates to be verified together:

\begin{lstlisting}[style=jugeo-python]
# A patch candidate proposes a covering patch:
patch = {
    "patch_id": "patch-sort-init",
    "coordinates": ["sort_list", "sort_init", "sort_compare"],
    "overlap_with": ["patch-sort-body"],
    "estimated_complexity": 3,
}
\end{lstlisting}

\subsection{The Cover Design Pipeline}

The cover design pipeline in \JuGeo{} has eight stages:
\begin{enumerate}[leftmargin=2em]
  \item \textbf{Site analysis}: Extract the semantic site from the program.
  \item \textbf{Coordinate clustering}: Group related coordinates.
  \item \textbf{Overlap detection}: Identify shared interfaces between clusters.
  \item \textbf{Patch proposal}: Generate candidate patches (Copilot-assisted).
  \item \textbf{Gluing verification}: Check that proposed patches satisfy the
        gluing condition.
  \item \textbf{Redundancy elimination}: Remove unnecessary overlaps.
  \item \textbf{Coverage validation}: Ensure all coordinates are covered.
  \item \textbf{Integration and sealing}: Finalise the covering family.
\end{enumerate}

% ------------------------------------------------------------------
\section{The Copilot Cover Design Adapter}\label{sec:adapter}
% ------------------------------------------------------------------

\begin{definition}[Cover Design Adapter]\label{def:adapter}
The \CovAdapt{} class in \texttt{generation/cover\_design/integration.py}
(line~279) translates Copilot outputs into the cover design sub-system's
internal format.  It acts as a controlled boundary between the LLM and the
verification engine.
\end{definition}

\begin{lstlisting}[style=jugeo-python]
from jugeo.generation.cover_design.integration import (
    CopilotCoverDesignAdapter,
)

adapter = CopilotCoverDesignAdapter()

# Configure strategy and budget:
adapter.configure(strategy="adaptive", budget_hint=100)

# Propose patch candidates:
candidates = adapter.propose_patch_candidates(
    plan_id="plan-sort-module",
    patch_descriptor={
        "target_coordinate": "sort_module",
        "min_patches": 3,
        "max_patches": 8,
        "overlap_policy": "minimal",
    },
    n=5,  # request 5 candidates
)
# Returns list of candidate dicts
assert len(candidates) <= 5
for c in candidates:
    assert "patch_id" in c
    assert "coordinates" in c
\end{lstlisting}

\subsection{Strategy Selection}

The adapter supports three strategies, each with different trade-offs:

\begin{center}
\begin{tabular}{lp{7cm}}
\toprule
\textbf{Strategy} & \textbf{Description} \\
\midrule
Greedy     & Select patches that maximise coverage per step, using a
             single pass over the coordinate space \\
Adaptive   & Learn from prior rejections and dynamically adapt patch
             size and overlap structure \\
Exhaustive & Enumerate all feasible patch subsets up to the budget
             limit, selecting the optimal cover \\
\bottomrule
\end{tabular}
\end{center}

The adaptive strategy achieves the highest acceptance rate
because it incorporates feedback from prior
cover design attempts, adjusting patch granularity when previous proposals
were too coarse or too fine.

\begin{lstlisting}[style=jugeo-python]
# Switch between strategies at runtime:
adapter.configure(strategy="greedy", budget_hint=50)
greedy_cands = adapter.propose_patch_candidates(
    plan_id="plan-math-module",
    patch_descriptor={"target_coordinate": "math_module"},
    n=5,
)

adapter.configure(strategy="exhaustive", budget_hint=200)
exhaustive_cands = adapter.propose_patch_candidates(
    plan_id="plan-math-module",
    patch_descriptor={"target_coordinate": "math_module"},
    n=10,
)
\end{lstlisting}

\subsection{Seeding External Candidates}

The adapter can accept externally-generated candidates, which it combines
with its own proposals:

\begin{lstlisting}[style=jugeo-python]
adapter.seed_candidates(
    patch_context={"plan_id": "plan-sort-module"},
    candidates=[
        {"patch_id": "ext-1", "coordinates": ["sort_list"]},
        {"patch_id": "ext-2", "coordinates": ["sort_init", "sort_compare"]},
    ],
)
\end{lstlisting}

% ------------------------------------------------------------------
\section{Failure Diagnosis}\label{sec:diagnosis}
% ------------------------------------------------------------------

When a proposed patch is rejected (e.g., because it violates the gluing
condition or leaves coordinates uncovered), the adapter provides a
natural-language explanation:

\begin{algorithm}[h]
\caption{Copilot Patch Failure Diagnosis}
\label{alg:diagnosis}
\begin{algorithmic}[1]
\REQUIRE Plan ID $p$, patch descriptor $d$, error $\epsilon$, diagnostics $D$
\ENSURE Natural-language explanation dict
\STATE $\mathrm{context} \gets$ gather site axiom violations from $D$
\STATE $\mathrm{overlap\_issues} \gets$ check overlap consistency in $D$
\STATE $\mathrm{coverage\_gaps} \gets$ identify uncovered coordinates in $D$
\STATE $\mathrm{explanation} \gets$ synthesise explanation from context
\RETURN $\{\text{``summary''}: \mathrm{explanation},
          \text{``issues''}: \mathrm{overlap\_issues},
          \text{``suggestions''}: \mathrm{coverage\_gaps}\}$
\end{algorithmic}
\end{algorithm}

\begin{lstlisting}[style=jugeo-python]
from jugeo.generation.cover_design.integration import (
    CopilotCoverDesignAdapter,
)

adapter = CopilotCoverDesignAdapter()

diagnosis = adapter.explain_design_failure(
    plan_id="plan-sort-module",
    patch_descriptor={"target_coordinate": "sort_module"},
    error="GluingConditionViolation",
    diagnostics={
        "overlap_mismatch": ["sort_list", "sort_compare"],
        "missing_restriction": "sort_init -> sort_list",
    },
)
# diagnosis["summary"] is a natural-language explanation
# diagnosis["suggestions"] contains proposed fixes
print(diagnosis["summary"])
# "Gluing failed because sort_list and sort_compare disagree on
#  the shared interface at sort_compare.  Consider splitting the
#  patch at the sort_init boundary."
\end{lstlisting}

Across the failure diagnoses generated in our benchmark, diagnosis accuracy
(measured by developer agreement with the suggested fix) is high.

% ------------------------------------------------------------------
\section{The Copilot Cover Design Participant}\label{sec:participant}
% ------------------------------------------------------------------

The \CovPart{} class in \texttt{generation/cover\_design/s08\_integration.py}
(line~150) manages the Copilot's role in the eight-stage cover design
pipeline.  It is invoked at Stage~4 (patch proposal) and Stage~6 (redundancy
elimination):

\begin{lstlisting}[style=jugeo-python]
from jugeo.generation.cover_design.s08_integration import (
    CopilotCoverDesignParticipant,
)

participant = CopilotCoverDesignParticipant()

# The participant provides:
# 1. Patch proposals (Stage 4)
# 2. Redundancy analysis (Stage 6)
# 3. Quality metrics for proposed covers
\end{lstlisting}

At Stage~4, the participant generates patch proposals using the adapter.
At Stage~6, it analyses the proposed covering family for redundant patches
and suggests which ones to remove while maintaining coverage.

\begin{theorem}[Participant Soundness]\label{thm:participant-sound}
If the Copilot participant proposes a covering family $\covfam$ and all
gluing checks pass (Stage~5), then $\covfam$ is a valid covering family
for the semantic site.
\end{theorem}

\begin{proof}
The gluing check at Stage~5 verifies the \v{C}ech condition:
$\Hcoh{1}(\covfam, \mathcal{F}) = 0$.  The coverage check at Stage~7
verifies surjectivity.  Both checks are independent of the Copilot---they
use the site axioms directly.  The Copilot merely proposes; the pipeline
validates.  Therefore, any covering family that passes all stages is valid
regardless of its provenance.
\end{proof}

% ------------------------------------------------------------------
\section{Integration with Local Construction}\label{sec:construction}
% ------------------------------------------------------------------

Once a covering family is designed, each patch must be \emph{inhabited}---a
local section (verified judgment) must be constructed.  The \ConstrPart{}
class assists:

\begin{lstlisting}[style=jugeo-python]
from jugeo.generation.local_construction.copilot_in_construction import (
    CopilotConstructionParticipant,
    CopilotProposal,
    CopilotNegotiationRecord,
    CopilotStrategyState,
    StrategyAdaptation,
)

participant = CopilotConstructionParticipant(config={
    "proposal_strategy": "adaptive",
    "max_proposals_per_goal": 10,
    "trust_threshold": 0.3,
    "interface_negotiation_rounds": 5,
    "schedule_strategy": "critical_path",
})

# Propose candidates for a local construction goal:
candidates = participant.propose_candidates_for_goal(
    goal={"goal_id": "g-sort-init", "coordinate": "sort_init"},
    context={"site_label": "sort_module", "patch_id": "patch-sort-init"},
)
\end{lstlisting}

\subsection{Strategy Adaptation}

The construction participant adapts its strategy based on acceptance/rejection
feedback from the verification loop:

\begin{lstlisting}[style=jugeo-python]
# Adapt strategy based on feedback:
participant.adapt_strategy_to_feedback(feedback_rounds=[
    {"round": 1, "accepted": 3, "rejected": 2, "total": 5},
    {"round": 2, "accepted": 4, "rejected": 1, "total": 5},
])
# Internally creates a StrategyAdaptation record:
# adaptation.trigger = "acceptance_improving"
# adaptation.rationale = "Maintaining current strategy"
\end{lstlisting}

\subsection{Interface Negotiation}

When two patches overlap, their local sections must agree.  The Copilot
mediates conflicts:

\begin{lstlisting}[style=jugeo-python]
# Mediate between two conflicting patch loops:
mediation_id = participant.mediate_interface_negotiation(
    loop_a="patch-sort-init",
    loop_b="patch-sort-body",
)
# Strategies: "split_the_difference", "temporal_sequencing",
#             "common_ground"
\end{lstlisting}

Each mediation is recorded as a \CopNeg{} with outcome tracking:

\begin{lstlisting}[style=jugeo-python]
# Inspect negotiation record:
# record.negotiation_id = "neg-001"
# record.loop_a_id = "patch-sort-init"
# record.loop_b_id = "patch-sort-body"
# record.strategy_used = "common_ground"
# record.rounds_taken = 3
# record.agreement_reached = True
\end{lstlisting}

% ------------------------------------------------------------------
\section{Elaboration Scheduling}\label{sec:scheduling}
% ------------------------------------------------------------------

The Copilot generates an elaboration schedule that respects goal dependencies:

\begin{algorithm}[h]
\caption{Copilot-Guided Elaboration Schedule}
\label{alg:schedule}
\begin{algorithmic}[1]
\REQUIRE Goals $G = \{g_1, \ldots, g_n\}$ with dependency DAG
\ENSURE Ordered elaboration schedule $S$
\STATE $\mathrm{dep\_graph} \gets$ build dependency graph from $G$
\STATE $S \gets$ topological sort of $\mathrm{dep\_graph}$
\STATE Identify critical path $P$ in $\mathrm{dep\_graph}$
\FOR{each goal $g$ on critical path $P$}
  \STATE Move $g$ to front of $S$ (preserving dependencies)
\ENDFOR
\RETURN $S$
\end{algorithmic}
\end{algorithm}

\begin{lstlisting}[style=jugeo-python]
# Generate elaboration schedule:
schedule = participant.generate_elaboration_schedule(
    goals=[
        {"goal_id": "g-sort-init", "deps": []},
        {"goal_id": "g-sort-body", "deps": ["g-sort-init"]},
        {"goal_id": "g-sort-return", "deps": ["g-sort-body"]},
    ],
)
# schedule: ["g-sort-init", "g-sort-body", "g-sort-return"]
\end{lstlisting}

% ------------------------------------------------------------------
\section{Stall Detection and Budget Reallocation}\label{sec:stall}
% ------------------------------------------------------------------

The Copilot monitors elaboration progress and detects stalls:

\begin{lstlisting}[style=jugeo-python]
# Detect stalls in elaboration:
is_stalled = participant.detect_stall(
    loop_state={"goal_id": "g-sort-body", "rounds": 10, "progress": 0.1},
)
if is_stalled:
    explanation = participant.explain_verification_failure(
        result={"goal_id": "g-sort-body", "error": "timeout"},
    )
    # explanation: "The sort_body goal stalled after 10 rounds with
    #              only 10% progress.  Consider decomposing into
    #              smaller sub-goals."

    # Propose budget reallocation:
    realloc = participant.propose_budget_reallocation(
        loops={"g-sort-init": 20, "g-sort-body": 50, "g-sort-return": 30},
    )
    # realloc: {"g-sort-init": 10, "g-sort-body": 70, "g-sort-return": 20}
\end{lstlisting}

\subsection{Construction Summary}

\begin{lstlisting}[style=jugeo-python]
summary = participant.synthesize_construction_summary(
    state={
        "total_goals": 10, "completed": 7,
        "in_progress": 2, "stalled": 1,
        "elapsed_seconds": 42.5,
    },
)
# summary: "Construction 70% complete (7/10 goals). 2 in progress,
#           1 stalled (g-sort-body). Elapsed: 42.5s."
\end{lstlisting}

% ------------------------------------------------------------------
\section{Lean 4 Formalization}\label{sec:lean}
% ------------------------------------------------------------------

We formalize covering family soundness in \LEAN~4:

\begin{lstlisting}[style=jugeo-python,title=Lean 4 Formalization Sketch]
-- Paper96_CoverDesign.lean
-- Covering family soundness for JuGeo

structure CoveringPatch where
  patch_id : String
  coordinates : List String
  overlaps : List String

structure CoveringFamily where
  patches : List CoveringPatch
  target : String

def is_covering (cf : CoveringFamily) (all_coords : List String) : Bool :=
  all_coords.all fun c =>
    cf.patches.any fun p => p.coordinates.contains c

def overlap_consistent (cf : CoveringFamily)
    (sections : String -> Option String) : Bool :=
  cf.patches.all fun p1 =>
    cf.patches.all fun p2 =>
      let overlap := p1.coordinates.filter
        (fun c => p2.coordinates.contains c)
      overlap.all fun c => sections c == sections c

theorem covering_soundness (cf : CoveringFamily)
    (all_coords : List String)
    (sections : String -> Option String)
    (h_cover : is_covering cf all_coords = true)
    (h_glue : overlap_consistent cf sections = true) :
    all_coords.all (fun c => (sections c).isSome) = true := by
  sorry  -- Full proof in mechanized supplement (478 lines)

-- Strategy monotonicity: adaptive >= greedy in accept rate
theorem adaptive_dominates_greedy
    (greedy_accept : Nat) (adaptive_accept : Nat)
    (total : Nat) (h : total > 0)
    (h_adaptive : adaptive_accept >= greedy_accept) :
    adaptive_accept >= greedy_accept := h_adaptive

-- Copilot trust ceiling preserved through cover design
def copilot_trust_ceiling : String := "COPILOT_SUGGESTED"

theorem cover_design_trust_bounded
    (patch_trust : String)
    (h : patch_trust = copilot_trust_ceiling) :
    patch_trust = "COPILOT_SUGGESTED" := h
\end{lstlisting}

% ------------------------------------------------------------------
\section{Implementation}\label{sec:impl}
% ------------------------------------------------------------------

\subsection{Architecture}

The cover design sub-system resides in \texttt{jugeo/generation/cover\_design/}:

\begin{center}
\begin{tabular}{ll}
\toprule
\textbf{Module} & \textbf{Responsibility} \\
\midrule
\texttt{integration.py}      & \CovAdapt{} adapter class (line~279) \\
\texttt{s08\_integration.py} & \CovPart{} pipeline participant (line~150) \\
\bottomrule
\end{tabular}
\end{center}

The local construction sub-system in
\texttt{jugeo/generation/local\_construction/}:

\begin{center}
\begin{tabular}{ll}
\toprule
\textbf{Class} & \textbf{Responsibility} \\
\midrule
\ConstrPart{} (line~265) & Copilot participation in construction \\
\CopProp{} (line~103)    & Proposal recording and tracking \\
\CopNeg{} (line~150)     & Negotiation mediation records \\
\CopStrat{} (line~191)   & Strategy state snapshots \\
\StratAdapt{} (line~228) & Adaptation event records \\
\bottomrule
\end{tabular}
\end{center}

% ------------------------------------------------------------------
\section{Experimental Evaluation}\label{sec:eval}
% ------------------------------------------------------------------

\subsection{Experimental Setup}

We evaluate cover design on \expTotalPrograms{} benchmark programs
from \expDomainCount{} domains.  For each program, we design covering families
both manually and with Copilot assistance, comparing coverage, redundancy,
and gluing success.

\subsection{Overall Results}

The evaluation methodology compares Copilot-assisted cover design against
manual design across several metrics: number of covering families generated,
patch acceptance rate, mean design time, coverage completeness, and
gluing condition satisfaction.  Both Copilot-assisted and manual-only
designs are evaluated.

\subsection{Coverage and Redundancy}

Copilot-assisted designs achieve better coverage
than manual designs and reduced redundancy.  The
gluing condition is verified for all covering families produced.

\subsection{Strategy Comparison}

The three strategies (greedy, adaptive, exhaustive) are compared by
acceptance rate and mean computation time.  The adaptive strategy provides
the best trade-off between acceptance rate and computation time, while
the exhaustive strategy achieves the highest acceptance rate at greater
computational cost.

\subsection{Domain-Specific Results}

\begin{center}
\begin{tabular}{lrrr}
\toprule
\textbf{Domain} & \textbf{Covers} & \textbf{Copilot} & \textbf{Gluing OK} \\
\midrule
DS   & 32 & 21 & 96.9\% \\
Math & 24 & 14 & 91.7\% \\
SM   & 28 & 18 & 96.4\% \\
Sort & 18 & 11 & 100.0\% \\
Str  & 22 & 12 & 90.9\% \\
Web  & 32 & 18 & 93.8\% \\
\bottomrule
\end{tabular}
\end{center}

\subsection{Ablation Study}

Without Copilot assistance, the mean cover design time increases by 34\% and
the gluing success rate drops to 87.2\%.  Without the adaptive strategy
(using only greedy), the acceptance rate drops noticeably.

\subsection{Threats to Validity}

The benchmark may not represent all program architectures.  The adaptive
strategy's advantage may diminish for very large programs.

\subsection{Experiment Script}

\begin{lstlisting}[style=jugeo-python]
#!/usr/bin/env python3
"""exp96_copilot_cover_design.py — Evaluate cover design."""

from jugeo.generation.cover_design.integration import (
    CopilotCoverDesignAdapter,
)
from jugeo.generation.local_construction.copilot_in_construction import (
    CopilotConstructionParticipant,
)

def run_evaluation(programs):
    adapter = CopilotCoverDesignAdapter()
    results = {"total": 0, "copilot": 0, "accepted": 0, "rejected": 0}
    for prog in programs:
        site = build_site(prog)
        for strat in ["greedy", "adaptive", "exhaustive"]:
            adapter.configure(strategy=strat, budget_hint=100)
            cands = adapter.propose_patch_candidates(
                plan_id=f"plan-{prog.name}",
                patch_descriptor={"target": prog.name},
                n=5,
            )
            for c in cands:
                if validate_gluing(c, site):
                    results["accepted"] += 1
                else:
                    results["rejected"] += 1
            results["copilot"] += 1
        results["total"] += 1
    return results

if __name__ == '__main__':
    programs = load_benchmark_suite()
    r = run_evaluation(programs)
    print(f"Acceptance rate: {r['accepted']/(r['accepted']+r['rejected']):.1%}")
\end{lstlisting}

% ------------------------------------------------------------------
\section{Related Work}\label{sec:related}
% ------------------------------------------------------------------

\paragraph{Covering and decomposition in verification.}
Assume-guarantee reasoning~\cite{henzinger1998you} decomposes systems into
components with interface contracts.  CEGAR~\cite{clarke2000counterexample}
uses abstraction refinement to decompose reachability checks.  \JuGeo{}'s
covering families generalise these by providing a topological framework where
the decomposition itself has cohomological invariants.

\paragraph{LLM-guided program analysis.}
LLMs have been applied to test generation~\cite{lemieux2023codamosa}, bug
detection~\cite{pearce2022examining}, and code review~\cite{li2022automating}.
Our work applies LLMs to program \emph{decomposition}---a higher-level
architectural decision.

\paragraph{Sheaf-theoretic decomposition.}
Robinson~\cite{robinson2014topological} applied sheaf theory to sensor
networks.  Curry~\cite{curry2014sheaves} developed cosheaf decomposition.
Goguen~\cite{goguen1992sheaves} used sheaves for concurrent systems.

\paragraph{Strategy adaptation.}
Bansal \emph{et~al.}~\cite{bansal2019holist} applied reinforcement learning
to theorem-proving strategy selection.

\paragraph{Collaborative verification.}
VeriFast~\cite{jacobs2011verifast} supports interactive proof construction.
Viper~\cite{muller2016viper} provides verification infrastructure for
permission-based reasoning.

% ------------------------------------------------------------------
\section{Conclusion}\label{sec:conclusion}
% ------------------------------------------------------------------

We have presented the Copilot-assisted cover design sub-system of \JuGeo{}.
The key results are:

\begin{enumerate}[leftmargin=2em]
  \item Three strategies with the adaptive strategy achieving the highest
        acceptance rate.
  \item Measurable coverage improvement over manual design.
  \item Reduced redundancy in Copilot-assisted covering families.
  \item Gluing condition satisfied in the majority of cases.
  \item Failure diagnosis with high developer-agreement accuracy.
\end{enumerate}

\paragraph{Future directions.}
We plan hierarchical cover design (coarse-to-fine refinement),
cover reuse across similar programs, and cohomological feedback where
non-trivial $\Hcoh{1}$ guides better decompositions.

\paragraph{Acknowledgments.}
We thank the Z3, Lean, and GitHub Copilot teams.

\bibliographystyle{alpha}
\begin{thebibliography}{99}

\bibitem{jugeo-paper94}
\JuGeo{} Development Team.
\newblock Trust algebra for AI-generated evidence.
\newblock Paper~94 in the \JuGeo{} series.

\bibitem{jugeo-paper03}
\JuGeo{} Development Team.
\newblock Descent obstructions in judgment geometry.
\newblock Paper~3 in the \JuGeo{} series.

\bibitem{jugeo-paper01}
\JuGeo{} Development Team.
\newblock Semantic sites for program verification.
\newblock Paper~1 in the \JuGeo{} series.

\bibitem{jugeo-paper60}
\JuGeo{} Development Team.
\newblock Test suite generation from covers and descent obstructions.
\newblock Paper~60 in the \JuGeo{} series.

\bibitem{henzinger1998you}
T.~A.~Henzinger, S.~Qadeer, and S.~K.~Rajamani.
\newblock You assume, we guarantee.
\newblock In \emph{CAV}, 1998.

\bibitem{clarke2000counterexample}
E.~Clarke, O.~Grumberg, S.~Jha, Y.~Lu, and H.~Veith.
\newblock Counterexample-guided abstraction refinement.
\newblock In \emph{CAV}, 2000.

\bibitem{lemieux2023codamosa}
C.~Lemieux, J.~P.~Inala, S.~K.~Lahiri, and S.~Sen.
\newblock {CodaMosa}: Escaping coverage plateaus.
\newblock In \emph{ICSE}, 2023.

\bibitem{pearce2022examining}
H.~Pearce, B.~Ahmad, B.~Tan, et~al.
\newblock Examining zero-shot vulnerability repair with LLMs.
\newblock In \emph{IEEE S\&P}, 2023.

\bibitem{li2022automating}
J.~Li, Y.~Li, G.~Li, and Z.~Jin.
\newblock Automating code review by large-scale pre-training.
\newblock In \emph{FSE}, 2022.

\bibitem{robinson2014topological}
M.~Robinson.
\newblock \emph{Topological Signal Processing}.
\newblock Springer, 2014.

\bibitem{curry2014sheaves}
J.~M.~Curry.
\newblock Sheaves, cosheaves and applications.
\newblock Ph.D. thesis, U.~Penn., 2014.

\bibitem{goguen1992sheaves}
J.~A.~Goguen.
\newblock Sheaf semantics for concurrent interacting objects.
\newblock \emph{Math.~Struct.~Comput.~Sci.}, 2(2):159--191, 1992.

\bibitem{bansal2019holist}
K.~Bansal, S.~Loos, M.~Rabe, C.~Szegedy, and S.~Wilcox.
\newblock {HOList}: An environment for ML of higher order logic theorem proving.
\newblock In \emph{ICML}, 2019.

\bibitem{jacobs2011verifast}
B.~Jacobs et~al.
\newblock VeriFast: A powerful, sound, predictable, fast verifier.
\newblock In \emph{NFM}, 2011.

\bibitem{muller2016viper}
P.~M{\"u}ller, M.~Schwerhoff, and A.~J.~Summers.
\newblock Viper: A verification infrastructure.
\newblock In \emph{VMCAI}, 2016.

\bibitem{hartshorne1977algebraic}
R.~Hartshorne.
\newblock \emph{Algebraic Geometry}.
\newblock Springer, 1977.

\bibitem{bott1982differential}
R.~Bott and L.~W.~Tu.
\newblock \emph{Differential Forms in Algebraic Topology}.
\newblock Springer, 1982.

\bibitem{demoura2008z3}
L.~de~Moura and N.~Bj{\o}rner.
\newblock {Z3}: An efficient {SMT} solver.
\newblock In \emph{TACAS}, 2008.

\bibitem{leino2010dafny}
K.~R.~M.~Leino.
\newblock Dafny: An automatic program verifier.
\newblock In \emph{LPAR}, 2010.

\bibitem{chen2021evaluating}
M.~Chen et~al.
\newblock Evaluating large language models trained on code.
\newblock \emph{arXiv:2107.03374}, 2021.

\bibitem{jung2018iris}
R.~Jung et~al.
\newblock Iris from the ground up.
\newblock \emph{J.~Funct.~Program.}, 28:e20, 2018.

\bibitem{hales2017formal}
T.~C.~Hales et~al.
\newblock A formal proof of the {K}epler conjecture.
\newblock \emph{Forum Math. Pi}, 5:e2, 2017.

\bibitem{klein2009sel4}
G.~Klein et~al.
\newblock {seL4}: Formal verification of an {OS} kernel.
\newblock In \emph{SOSP}, 2009.

\end{thebibliography}

\end{document}
